% Problem setup and geometric motivation
\section{Problem setup and geometric motivation}
\label{sec:notation}

Let $\mathrm{enc}_\ell : \Sigma^* \to \mathbb{R}^d$ be the hidden state at layer $\ell$ of a language model over vocabulary $\Sigma$. For each sample $i$ we have a representation $\mathbf{x}_i = \mathrm{enc}_\ell(s_i) \in \mathbb{R}^d$, a binary concept label $y_i \in \{0,1\}$, and an optional control label $z_i$. The dataset is $\mathcal{D} = \{(\mathbf{x}_i, y_i, z_i)\}_{i=1}^N$ with class subsets $\mathcal{D}_c = \{i : y_i = c\}$. A concept scorer $f_\theta : \mathbb{R}^d \to \mathbb{R}$ is a 2-layer MLP with input gradient $\nabla f_\theta(\mathbf{x}_i)$.

\paragraph{Local manifold hypothesis.} We adopt the modeling assumption that the representations $\{\mathbf{x}_i\}$ lie near a locally low-dimensional support $\mathcal{M} \subset \mathbb{R}^d$. Its operational consequence is that we can approximate the neighborhood of any $\mathbf{x} \in \mathcal{M}$ by its tangent space $T_{\mathbf{x}}\mathcal{M}$ \citep{tu-2011-manifolds} and apply linear-algebraic operations (SVD, projection, anisotropic scaling) inside it. The tangent space is not observed directly; we estimate it from nearby representations by local SVD on the secant matrix (§\ref{sec:method}). For an ambient vector $\mathbf{g} \in \mathbb{R}^d$, the decomposition $\mathbf{g} = \mathbf{g}^{\top} + \mathbf{g}^{\perp}$ with $\mathbf{g}^{\top} \in T_{\mathbf{x}}\mathcal{M}$ separates the part of $\mathbf{g}$ supported by the local sample structure from the part orthogonal to it. \TanGCE{} acts on the tangent component only; \AmbGCE{} retains both. Formal coordinate-chart definitions and the differential $d(\varphi^{-1})_0$ identification with $T_{\mathbf{x}}\mathcal{M}$ are deferred to App.~\ref{sec:appendix_geometry}.

\paragraph{Recoverability.} Concept removal is judged only after training a \emph{fresh} probe on the edited representations \citep{elazar-goldberg-2018-adversarial}: erasure succeeds if a newly trained nonlinear classifier fails to recover the concept from $\tilde{\mathbf{x}}$, not merely if the old scorer is fooled.
